\chapter{Datensplit und Training} 
\label{chap:datensplit_und_training}

\section{Anforderungen an die Datenbasis}

Um eine konsistente, validierte und reproduzierbare Datenbasis zu haben, wurde ein idealer Datensatz erstellt.
Der ideale Datensatz setzt sich aus den bestgeeigneten Artikeln und Makrokandidaten nach der Datenbereinigung zusammen \parencite{srcIdealDataset,srcStep913}.
Dabei wurden die Anforderungen an den idealen Datensatz in Tabelle \ref{tab:anforderungen_datenbasis} zusammengefasst.

\begin{table}[tbp]
    \centering
    \setlength{\tabcolsep}{6pt}
    \renewcommand{\arraystretch}{1.08}
    \caption{Anforderungen an den idealen Datensatz \parencite{srcIdealDataset,srcStep913,srcStep921}}
    \label{tab:anforderungen_datenbasis}
    \begin{tabularx}{\textwidth}{>{\raggedright\arraybackslash}p{4.2cm} >{\raggedright\arraybackslash}X}
        \toprule
        \textbf{Anforderung} & \textbf{Beschreibung und Zielsetzung} \\
        \midrule
        Reale Datenbasis & Der Datensatz enthält ausschließlich tatsächlich vorhandene EPLAN-Artikel aus der gesamten Datenbasis. \\
        Eindeutiges Zielmakro & Jeder Artikel muss eindeutig einem Zielmakro (\texttt{target\_macro}) zugeordnet sein. \\
        Vorhandene Ground Truth & Für jedes Zielmakro muss eine passende reale \texttt{.ema}-Makrodatei existieren. \\
        Vollständige technische Merkmale & Die Feature-Vektoren müssen vorhanden und mit technischen Informationen gefüllt sein. \\
        Vollständige Grunddaten & Jeder Datensatz benötigt eine ID, ein Zielmakro, eine Makrofamilie und einen Fingerprint. \\
        Keine Duplikate & Artikel-IDs dürfen nicht mehrfach vorkommen, damit kein Artikel versehentlich doppelt trainiert wird. \\
        Keine widersprüchlichen Fingerprints & Ein Fingerprint darf nicht zu mehreren unterschiedlichen Zielmakros gehören. Andernfalls wäre die Zuordnung technisch nicht eindeutig. \\
        Ausgeglichener Target-Support & Jedes Zielmakro sollte ungefähr gleich häufig vertreten sein. Dadurch lernt das Modell nicht überwiegend die am häufigsten vorkommenden Makros. \\
        Vergleichbare Makrofamilien & Der Datensatz sollte mehrere Zielmakros derselben Familie enthalten. So kann geprüft werden, ob das Modell auch zwischen ähnlichen Makros unterscheiden kann. \\
        Leakage-freie und reproduzierbare Aufteilung & Die Daten werden mit einem festen Zufallswert (\texttt{RANDOM\_SEED} = 42) in 70\,\% Training, 10\,\% Validierung und 20\,\% Test aufgeteilt. Artikel mit zusammengehörigem Fingerprint bleiben dabei im selben Teil des Datensatzes \parencite{srcStep921}.\\
        \bottomrule
    \end{tabularx}
\end{table}

\section{Datensplit und Trainingsbasis}

Nach der Datenaufbereitung, der Qualitätsprüfung der Daten sowie der Erstellung des idealen Datensatzes erfolgt die Aufteilung in Trainings-, Validierungs- und Testdaten \parencite{srcStep911,srcStep913,srcIdealDataset,srcStep921}.
Ziel ist es eine klare Trennung zwischen Modelltraining, Hyperparameteroptimierung und Modellbewertung zu schaffen. 

\begin{align}
    \text{Train} = 70\,\%, \qquad \text{Validation} = 10\,\%, \qquad \text{Test} = 20\,\%
\end{align}

Um reproduzierbare Ergebnisse mit einem erstellten Datensplit zu gewährleisten, wird ein fester Zufallsseed verwendet, der sicherstellt, dass die Aufteilung in Trainings-, Validierungs- und Testdaten bei jeder Ausführung identisch bleibt \parencite{srcStep921}.

\begin{align}
    \text{RANDOM\_SEED} = 42
\end{align}

Dabei wird keine reine Zufallsaufteilung vorgenommen, sondern fachliche Zusammenhänge der Daten berücksichtigt, um eine konsistente Verteilung über die Anfragegruppen sicherzustellen. 
Dies verhindert, dass eng verwandte Beispiele in unterschiedlichen Datensplits landen und somit die Glaubwürdigkeit der Modellbewertung negativ beeinflussen. Durch die unterschiedlichen Seed-Werte lässt sich dann auch ein Mittelwert über mehrere Splits ableiten \parencite{srcStep921}. Abbildung \ref{fig:datensplit_schema} fasst diese Aufteilung schematisch zusammen.

\begin{figure}[tbp]
    \centering
    \begin{tikzpicture}[
        font=\small,
        bar/.style={draw=blue!60!black, thick, rounded corners=4pt, align=center, minimum height=2.2cm, inner sep=5pt},
        seed/.style={draw=gray!50!black, rounded corners=4pt, fill=gray!8, align=center, minimum height=1.1cm, inner sep=5pt}
    ]
        \node[bar, fill=blue!15, text width=5.6cm] (train) at (0,0) {\textbf{Trainings-Split (70\,\%)}\\[1pt]\footnotesize Modelltraining \& Feature-Extraktion};
        \node[bar, fill=orange!20, text width=3.6cm, right=3pt of train] (val) {\textbf{Validierungs-Split (10\,\%)}\\[1pt]\footnotesize Hyperparameter-\\optimierung};
        \node[bar, fill=green!18, text width=3.2cm, right=3pt of val] (test) {\textbf{Test-Split (20\,\%)}\\[1pt]\footnotesize Evaluation};

        \node[seed, below=0.3cm of train.south west, anchor=north west, text width=13.3cm] (seedbox) {\footnotesize \textbf{Reproduzierbarkeit:} \texttt{RANDOM\_SEED = 42} $\mid$ Fingerprint-gruppierter Split (Leakage-Schutz)};
    \end{tikzpicture}
    \caption{Schematische Aufteilung der Trainings-, Validierungs- und Testdatenbasis.}
    \label{fig:datensplit_schema}
\end{figure}

\section{Aufbau der Retrievalsignale}

Vor dem Modelltraining werden für jede Anfrage unterschiedliche Retrievalsignale berechnet (siehe Kapitel \ref{chap:analyse_der_datenstrukturen}) \parencite{srcStep640}.
Diese dienen als erste Ähnlichkeitsabschätzung zwischen Anfrage und Kandidat.

Die theoretischen Grundlagen der verwendeten Verfahren wurden bereits im Kapitel \ref{chap:stand_der_technik} beschrieben.
Die Retrievalsignale dienen innerhalb der in Abbildung \ref{fig:architektur} dargestellten Pipeline-Architektur nicht als eigenständige Rankingstufen, sondern als zusätzliche Merkmale für das Learning-to-Rank-Modell \parencite{srcStep640,srcStep931}.

Die folgende Tabelle \ref{tab:retrieval_signale} zeigt die verwendeten Retrievalsignale sowie deren Funktion.

\begin{table}[tbp]
    \centering
    \setlength{\tabcolsep}{6pt}
    \renewcommand{\arraystretch}{1.1}
    \caption{Verwendete Retrievalsignale} 
    \label{tab:retrieval_signale} 
    \begin{tabularx}{\textwidth}{>{\raggedright\arraybackslash}p{3.5cm} >{\raggedright\arraybackslash}X}
        \toprule
        \textbf{Signal} & \textbf{Funktion} \\
        \midrule
        BM25 & Lexikalische Ähnlichkeit zwischen Anfrage und Makrokandidatenprofil \\
        Exact Match & Exakte technische Übereinstimmungen strukturierter Merkmale \\
        RRF & Fusion der Retrievalrankings zu einem kombinierten Retrievalsignal \\
        \bottomrule 
    \end{tabularx}
\end{table}


Die verwendeten Parameter der Retrievalkomponenten wurden wie folgt festgelegt \parencite{srcStep640}:

\begin{align}
    \text{BM25}_{k1} &= 1.2 \\
    \text{BM25}_{b}  &= 0.50 \\
    \text{RRF}_{k}   &= 30
\end{align}

Dabei steuert $k_1$ wie stark wiederholte Treffer eines Begriffes berücksichtigt werden. Hier wurden 1,2 gewählt, da sich wiederholende lexikalische Begriffe zwar relevant sein können,
aber nicht automatisch bedeuten, das ein EPLAN-Artikel \parencite{EPLANArtikelverwaltung} für ein EPLAN-Fenstermakro \parencite{EPLAN-Makros} relevanter ist.

Der Parameter $b$ steuert die Normalisierung der Dokumentenlänge. Dieser liegt zwischen 0 und 1 und wurde hier auf 0,50 gesetzt. Dadurch wird berücksichtigt, dass längere Kandidatenprofile mehr Treffer enthalten können,
diese aber nicht aufgrund der Länge einen übermäßigen Vorteil erhalten \parencite{srcStep640}.

$\text{RRF}_k$ steuert den Einfluss der Spitzenpositionen. Bei kleinem $k$ dominiert Rang 1 stark, während bei größerem $k$ auch niedrigere Ränge stärker berücksichtigt werden.
60 ist die übliche Wahl. Hier wurde 30 verwendet, um die Spitzenpositionen nicht zu stark zu reduzieren und damit dem Ranking auf Top 1 mehr Abstand zu verleihen \parencite{srcStep640}.

Die Fusion der Retrievalsignale erfolgt über die Gewichtung der Reciprocal Rank Fusion (RRF) \parencite{srcStep640}:

\begin{align}
    w_{\text{BM25}} = 0.30 \qquad \text{und} \qquad w_{\text{Exact}} = 0.70
\end{align}

Die höhere Gewichtung von Exact Match gegenüber BM25 spiegelt die Priorität technischer Merkmale innerhalb der Maschinen- und Anlagenbau-Domäne wider. 


\section{Training des Learning-to-Rank-Modells}

Das zentrale Element der Architektur ist das Learning-to-Rank-Modell. Hierfür wird der \texttt{LGBMRanker} der LightGBM-Bibliothek verwendet \parencite{srcStep931}.
Dieser basiert auf dem in Kapitel \ref{chap:stand_der_technik} beschriebenen LambdaRank-/LambdaMART-Prinzip.

Das Modell wird mit den Parametern 

\begin{align} 
    \text{objective} = \text{lambdarank} \qquad \text{und} \qquad \text{metric} = \text{ndcg}
\end{align} 

trainiert \parencite{srcStep931}.


Wie in Kapitel \ref{chap:analyse_der_datenstrukturen} beschrieben, werden Kandidaten einer Anfrage als gemeinsame Rankinggruppe betrachtet.
Wiederholend entsteht für jede Anfrage eine Gruppe:

\begin{align} 
    G_q = \{(q,c_1),(q,c_2),\dots,(q,c_n)\} 
\end{align} 

mit der Anfrage $q$ und den zugehörigen Makrokandidaten $c_i$.

Dadurch erlernt das Modell nicht die allgemeine Zugehörigkeit eines einzelnen Kandidaten zu einer Klasse. Stattdessen wird die optimale Kandidatenreihenfolge innerhalb derselben Anfrage gelernt.
Ziel dabei ist es, dass der relevanteste Kandidat möglichst weit oben in der Rangliste erscheint.

Die Eingaben für das Modell sind die im Kapitel \ref{chap:analyse_der_datenstrukturen} beschriebenen Pairwise-Merkmale.
Das Training erfolgt dabei ausschließlich auf den Trainingsdaten, während die Optimierung der Hyperparameter über die Validierungsdaten erfolgt \parencite{srcStep931}.

\section{Hyperparameteroptimierung}

Die Leistungsfähigkeit eines Learning-to-Rank-Modells hängt von der Wahl der Hyperparameter ab. Deswegen erfolgt die Modellkonfiguration nicht manuell, sondern durch eine systematische Hyperparameteroptimierung.
Dafür wird hier eine Grid-Search-basierte Strategie verwendet. Bei dieser werden verschiedene Kombinationen von Hyperparametern getestet und anhand der Ergebnisse der Validierung verglichen \parencite{srcStep931}.

Die folgende Tabelle \ref{tab:hyperparameter} zeigt die untersuchten Hyperparameterkombinationen:

\begin{table}[tbp] 
    \centering 
    \setlength{\tabcolsep}{8pt}
    \renewcommand{\arraystretch}{1.1}
    \caption{Verwendete Hyperparameterkombinationen \parencite{srcStep931}}
    \label{tab:hyperparameter} 
    \begin{tabularx}{0.85\textwidth}{>{\raggedright\arraybackslash}p{5.5cm} >{\raggedright\arraybackslash}X} 
        \toprule 
        \textbf{Parameter} & \textbf{Wertebereich} \\ 
        \midrule 
        \texttt{n\_estimators} & 300, 500, 700 \\ 
        \texttt{learning\_rate} & 0.05, 0.04, 0.035 \\ 
        \texttt{num\_leaves} & 31, 63 \\ 
        \texttt{min\_child\_samples} & 10, 20, 50 \\ 
        \texttt{lambdarank\_\allowbreak truncation\_level} & 10, 20 \\ 
        \bottomrule 
    \end{tabularx} 
\end{table}

Dabei sind die Hyperparameter wie folgt zu verstehen:

\begin{itemize}
    \item \textbf{\texttt{n\_estimators}:} Anzahl der aufgebauten Entscheidungsbäume im Modell. Mehr Bäume = komplexere Zusammenhänge, aber mehr Rechenaufwand.
    \item \textbf{\texttt{learning\_rate}:} Lernrate pro Trainingsschritt. Kleinere Werte führen zu vorsichtigerem, aber stabileren Lernen.
    \item \textbf{\texttt{num\_leaves}:} Maximale Anzahl der Blätter pro Baum. Mehr Blätter = komplexere Bäume, aber höheres Risiko für Overfitting.
    \item \textbf{\texttt{min\_child\_samples}:} Minimale Anzahl an Trainingsbeispielen in einem Blatt. Kleinere Werte = mehr Detailtreue pro Blatt, größere Werte = bessere Regularisierung (verhindert Overfitting).
    \item \textbf{\texttt{lambdarank\_\allowbreak truncation\_level}:} Begrenzung des Ranglistenbereichs, der durch LambdaRank optimiert werden soll (Top 10, Top 20).
\end{itemize}

Für jede Parameterkombination wird ein vollständiger Trainingsdurchlauf durchgeführt. Die Bewertung erfolgt anschließend auf dem Validierungsdatensatz.
Das finale Modell wird auf Basis der Rankingqualität anhand der Validierungsdaten bestimmt \parencite{srcStep931}.

Die verwendeten Bewertungsmetriken werden später in Kapitel \ref{chap:evaluierungskonzept_und_ergebnisse} ausführlich beschrieben.

\section{Modellartefakte und Reproduzierbarkeit}

Ein wichtiges Ziel dieser Trainingspipeline besteht darin, die Trainingsabläufe reproduzierbar zu gestalten. Deswegen werden nach jedem Trainingslauf alle relevanten Artefakte gespeichert und dokumentiert \parencite{srcStep931}. Abbildung \ref{fig:artefakte_struktur} zeigt die zugehörige Ablagestruktur.

\begin{figure}[tbp]
    \centering
    \begin{tikzpicture}[
        font=\small,
        dirbox/.style={draw=blue!60!black, thick, rounded corners, fill=blue!5, text width=4.8cm, align=left, inner sep=6pt},
        artbox/.style={draw=teal!60!black, thick, rounded corners, fill=teal!5, text width=7.0cm, align=left, inner sep=6pt}
    ]
        \node[dirbox] (dir) {
            \textbf{Ablagestruktur im Laufverzeichnis:}\\[4pt]
            \ttfamily\small
            run\_directory/\\
            $\text{\textbar}$-- model/\\
            $\text{\textbar}$-- training/\\
            $\text{\textbar}$-- evaluation/\\
            $\text{\textbar}$-- summary/\\
            $\text{\textbackslash}$-- error\_analysis/
        };
        \node[artbox, right=0.5cm of dir] (art) {
            \textbf{Gespeicherte Modellartefakte:}\\[4pt]
            \small
            \begin{itemize}[leftmargin=12pt, itemsep=2pt, topsep=0pt]
                \item das trainierte LightGBM-Modell,
                \item die verwendeten Hyperparameter,
                \item die erzielten Validierungsmetriken,
                \item die verwendeten Featuredefinitionen,
                \item die Makrokandidatenprofile,
                \item Konfigurations- und Laufmetadaten.
            \end{itemize}
        };
    \end{tikzpicture}
    \caption{Strukturierte Ablage der Modellartefakte im Laufverzeichnis.}
    \label{fig:artefakte_struktur}
\end{figure}

Dadurch lässt sich sicherstellen, dass jeder Trainingsablauf zu einem späteren Zeitpunkt vollständig nachvollzogen und reproduziert werden kann \parencite{srcStep931}.


\section{Pipeline-Orchestrierung}

Die Ausführung der gesamten Pipeline erfolgt über eine zentrale Komponente, die die Orchestrierung übernimmt. 
Dadurch wird sichergestellt, dass alle Verarbeitungsschritte in der richtigen Reihenfolge ausgeführt werden \parencite{srcStep936}.

Die folgende Reihenfolge der Schritte stellt den Ablauf der Pipeline dar. Diese können auch aus dem Architekturdiagramm in Abbildung \ref{fig:architektur} entnommen werden:

\begin{enumerate} 
    \item Datenaufbereitung 
    \item Qualitätsprüfung und Datenvalidierung 
    \item Datensplit   
    \item Berechnung der Retrievalsignale 
    \item Erstellung der Pairwise-Features 
    \item Training des Rankingmodells 
    \item Speicherung der Modellartefakte 
    \item Evaluation des trainierten Modells 
\end{enumerate}

Jeder Verarbeitungsschritt verwendet dabei die Artefakte des vorherigen Schrittes als Eingabe und erzeugt selbst Artefakte für den nächsten Schritt als Ausgabe. 
Dadurch wird ein konsistenter Datenfluss von den Roh-Daten bis zum trainierten Modell sichergestellt \parencite{srcStep936,srcPipeline}.